\par\noindent\SourceLeaf{54}%
Sit in hac expressione $(2q-1) = r$, habebimus hac substitutione ipsius $\cos r\varphi$ (existente $r$ impari) coefficientem ita constitutum, ut sit eius terminus initialis

\[
= \frac{1}{2^{r-1}} {}^{\nu}\mathfrak{B}^{(r)} n^{r},
\]

\noindent terminus $h^{\text{tus}}$

\[
= \frac{1}{2^{r+2h-1}} {}^{r+2h}\mathfrak{B}^{(r+h)}  .  {}^{\nu}\mathfrak{B}^{(r+2h)} n^{r+2h}.
\]

\noindent Erat porro coefficientis ad $\cos 2q\varphi$ pertinentis terminus initialis

\[
\frac{1}{2^{2q-1}} {}^{\nu}\mathfrak{B}^{(2q)} n^{2q}, \qquad h^{\text{tus}} = \frac{1}{2^{2q+2h-1}} {}^{2q+2h}\mathfrak{B}^{(2q+h)}  .  {}^{\nu}\mathfrak{B}^{(2q+2h)} n^{2q+2h}.
\]

\noindent Sit eadem ratione $2q = r$, veniet hac substitutione pro coefficiente ipsius $\cos r\varphi$ (existente $r$ pari) expressio ita comparata, ut sit eius terminus initialis

\[
= \frac{1}{2^{r-1}} {}^{\nu}\mathfrak{B}^{(r)} n^{r},
\]

\noindent generalis $h^{\text{tus}}$

\[
= \frac{1}{2^{r+2h-1}} {}^{r+2h}\mathfrak{B}^{(r+h)}  .  {}^{\nu}\mathfrak{B}^{(r+2h)} n^{r+2h}.
\]

\noindent Identitas expressionum, quas, sive par sit sive impar assumtus $r$, pro coefficiente ipsius $\cos r\varphi$ adipiscimur, identitatem legis arguit, quam igitur, pro omnibus ipsius $(1+n\cos\varphi)^{\nu}$ terminis unam eandemque, concinnius finitima hac formula designabimus.

Est igitur quantitas nostra $(1+n\cos\varphi)^{\nu}$ seriei aequalis, quae sic procedit, ut sit

\[
A + A^{(1)}\cos\varphi + A^{(2)}\cos 2\varphi + .\,.\,.\,. + A^{(r)}\cos r\varphi + .\,.\,.
\]

\noindent ea coefficientium lege, ut initialis $A$ aequetur seriei, cuius terminus initialis $=1$, generalis $h^{\text{tus}}$

\[
= \frac{1}{2^{2h}} {}^{2h}\mathfrak{B}^{(h)}  .  {}^{\nu}\mathfrak{B}^{(2h)} n^{2h};
\]

\noindent cuius coefficiens indefinitus $r^{\text{tus}}$ $A^{(r)}$ contineatur serie, cuius terminus initialis

\[
= \frac{1}{2^{r-1}} {}^{\nu}\mathfrak{B}^{(r)} n^{r};
\]

\noindent coefficiens termini in eadem serie $h^{\text{ti}}$

\[
= \frac{1}{2^{r+2h-1}} {}^{\nu}\mathfrak{B}^{(r+2h)}  .  {}^{r+2h}\mathfrak{B}^{(r+h)} n^{r+2h}.
\]
